\documentclass{article}

\usepackage{geometry}
\geometry{a4paper}

\usepackage[round,authoryear]{natbib}
\bibliographystyle{plainnat}
\usepackage[hidelinks]{hyperref}
\usepackage{authblk}
\usepackage{tabularray}
\usepackage{graphicx}
\usepackage{amssymb}
\usepackage{xcolor}

% Make affiliations run inline
\renewcommand\Affilfont{\small}
\renewcommand\Authfont{\normalsize}

\makeatletter
\renewcommand\AB@affilsepx{;\hspace{0.6em}\protect\Affilfont}
\makeatother

\renewcommand{\thetable}{S\arabic{table}}
\renewcommand{\thefigure}{S\arabic{figure}}
\renewcommand{\thesubsection}{\arabic{subsection}}

\begin{document}

\title{Population-scale Ancestral Recombination Graphs with tskit 1.0\\
\large Supplementary Information}

\input{authors.tex}

\date{}
\maketitle

\subsection*{Published tools using tskit}
We performed a search for published software using tskit.
To be included in this list, the software must be described in a publication
as a named tool intended for reuse and must use tskit.
This list is intentionally conservative, and does not include
(for example) publications depending on bespoke analysis
pipelines using tskit.
For each tool in Table S1, we list: its name (with a clickable
hyperlink to its source code repository);
the main implementation
language(s) of the tool based on GitHub summaries;
the DOI (with clickable hyperlink) of the publication in which this
tool was first described (or, in the case of SLiM, the first
publication in which tskit-based features were described);
and the year of publication and number of citations to this
publication according to data from OpenAlex
(\url{https://openalex.org}) retrieved on 2026-01-23.

\input{tools_table.tex}

\subsection*{Tskit functionality and benchmarks}

In this section we review tskit's features and some available benchmarks
that compare it to other approaches.
Tskit is the result of more than a decade of sustained development,
and its distinctive contribution is to combine a general-purpose ARG
representation with efficient primitive operations for traversal,
editing, and analysis.
At its foundation is the succinct tree sequence encoding,
first introduced as part of the msprime coalescent
simulator~\citep{kelleher2016efficient}, together with the
sequential tree-construction algorithm for iterating over
trees along the genome.
The key algorithmic result is that
the cost of iterating over the local trees along the genome scales
with the number of trees and the edge differences between neighbouring
trees, rather than with the sample size.
The practical consequence was demonstrated in a benchmark in~\citet{kelleher2016efficient} 
which showed that using the succinct tree sequence format and sequential iteration was more
than a million times faster than reading the same trees individually
from Newick using the standard libraries available at the time, for a
simulation of 100,000 samples, with storage requirements reduced by
roughly four orders of magnitude.
The data model was subsequently generalised to enable forward-time
simulation in~\citet{kelleher2018efficient}. This publication also
introduced the columnar tables API and was the first paper in which
the tskit library itself appeared as a distinct component.
That work introduced the ``simplify'' operation, which restricts an ARG
to the history ancestral to a specified set of samples, showing that
its time complexity is linear in the number of input edges.
ARG simplification is a fundamental operation with many
applications~\citep{wong2024general}, as well as being the primitive
operation making forward-time simulation of ARGs
tractable~\citep{haller2018tree}.

Building on this foundation, tskit has accumulated a broad set of
reusable analytical primitives.
Motivated by the need to summarise biobank-scale inferred ARGs with
uncalibrated node times, \citet{kelleher2019inferring} introduced the
genealogical nearest neighbours (GNN) statistic, and demonstrated that
the tree-sequence encoding compresses inferred biobank-scale data by
orders of magnitude relative to VCF.
\citet{ralph2020efficiently} then introduced a general API for
population-genetic statistics built around a duality theorem relating
site-based and branch-based statistics. The benchmarks in that paper
show speedups of roughly three orders of magnitude, together with a
$\sim 200$-fold reduction in memory footprint, against genotype-matrix
implementations at a sample size of one million haplotypes.
The msprime~1.0 release~\citep{baumdicker2022efficient} provided
further illustrations of the flexibility of the data model, including
simulations of gene conversion and complete ARGs.
Individual and pedigree abstractions were added to the tskit data
model by \citet{andersontrocme2023genes}, extending what a node or
individual in an ARG can represent.
\citet{tsambos2023link} introduced the ``link\_ancestors'' operation,
a new primitive for extracting local-ancestry tracts from an ARG.
Development of the core API has continued through 2024--2026.
\citet{fritze2026forest} added the ``extend\_haplotypes'' function to
tskit, which restores unary-node spans lost to simplification or
inference and roughly halves the number of edges in typical ARGs, with
corresponding reductions in statistic-computation time. \citet{lehmann2026on} 
introduced operations to calculate ``branch'' genetic relatedness matrices 
(GRMs), which encode pairwise relatedness via shared branch area in an ARG, and their 
vector products directly from the ARG. As described in \citet{lehmann2026on},
for one million simulated diploids on human chromosome~21, the 
branch GRM--vector product completes in under 18~seconds, and the associated
randomised PCA scales to $2^{20}$ samples, a scale at which genotype-matrix PCA
would run out of memory on most machines.
\citet{lee2026genetic} apply this last operation to fitting linear mixed models,
and compare the branch GRM-vector product to the corresponding approximate (Monte Carlo)
implementation in ARGneedle-lib~\citep{zhang2023biobank}, showing that tskit is
one to two orders of magnitude faster despite using an exact algorithm.

A second defining feature of tskit is composability across different
stages of analysis.
Because simulation, inference, dating, and downstream analysis all
operate on the same underlying representation, methods developed in
one setting can be reused directly in another.
The pyslim library~\citep{gopalan2025bridging} bridges SLiM and tskit
at the Python level, operationalising these hybrid workflows together
using tskit's ``union'' operation for stitching together tree
sequences from independent simulations; this combination has enabled
whole-chromosome forward simulations of the entire 10-million-year
history of the great apes~\citep{rodrigues2024shared} and
whole-genome human simulations under realistic demographic and
selection models~\citep{haller2025simHumanity}.
\citet{wohns2022unified} introduced the tsdate method for estimating
ages of ancestral nodes from an inferred ARG under a mutational-clock
prior, and used it to combine 3{,}601 modern and eight high-coverage
ancient human genomes into a single unified genealogy.
Alongside these core additions, tskit's exposure to Numba~\citep{lam2015numba} 
of vectors allowing structured tree traversal
has allowed Python-level user code to JIT-compile ARG traversals at close to C speed,
enabling downstream packages such as
tstrait~\citep{tagami2024tstrait} and
tsbrowse~\citep{karthikeyan2025tsbrowse}
to perform biobank-scale computations in pure Python.

Real-data applications have further stressed the data model outside
its original coalescent-simulation setting: an ARG for 2.48~million
SARS-CoV-2 genomes inferred with sc2ts occupies only 32\,MiB and
loads into memory in about half a
second~\citep{zhan2025pandemic}; and recent work has dated contiguous
ARGs spanning 14\% of the human genome for 47,535 Genomics England
participants, estimating ages for 23.2 million
variants~\citep{pope2026tracing}.
More broadly, Table~S1 shows that tskit now underpins a substantial
ecosystem of downstream tools spanning simulation, ARG inference,
population-genetic and statistical-genetic inference, analysis, and
visualisation.

Tskit~1.0 is the point at which this decade-long accumulation of
representation, algorithms, and reusable operations is brought
together and stabilised under a single, guaranteed API.
Tskit is a software library and stable data model for representing
ARGs using the succinct tree sequence encoding, and for performing
core traversal, editing, and statistical operations on them.
To situate tskit relative to existing software, we contrast its
functionality with three Python-accessible libraries from adjacent
niches.
We compare with ARGneedle-lib~\citep{zhang2023biobank}, which infers
and analyses ARGs from biobank-scale data, because it is the only
other currently available Python-accessible software library that
provides a lossless ARG representation.
We compare with BTE (\url{https://github.com/jmcbroome/BTE}), the
Python interface to matUtils~\citep{mcbroome2021daily}, a toolkit for
manipulating and analysing very large mutation-annotated trees of
SARS-CoV-2, to illustrate the overlap between tskit and a
virus-specialised tree toolkit.
Finally, we compare with DendroPy~\citep{sukumaran2010dendropy}, a
mature and widely used phylogenetic scripting library with a rich
suite of tree operations.
Together, these comparisons illustrate both the breadth of general
operations supported by tskit and the distinctive features that arise
from its underlying ARG data model.
It is important to note that all these libraries have different goals than tskit, and
we are comparing their functionality in terms of the core ARG operations implemented by tskit,
rather than the main use cases addressed by the authors of these libraries.

Table S2 lists properties and operations that tskit supports, grouped
into nine broad categories.
For each operation we mark whether each comparison library offers a
comparable feature:
\begin{itemize}
\item[\checkmark{}] denotes full support, where the feature is
documented and accessible in the Python API, comparable in scope to
the tskit equivalent;
\item[$\circ$] denotes partial or restricted support, where the
feature is accessible but narrower in scope than the tskit
equivalent, is available only indirectly via existing operations, or
is not represented as a first-class part of the library's data model;
\item[] a blank cell indicates that no comparable feature is
available through the Python API.
\end{itemize}
Our comparison is intended to assess reusable, documented library
functionality rather than undocumented internal implementation
details.
Although ARGneedle-lib and matUtils are implemented in C++, neither
provides a documented low-level API clearly intended for downstream reuse, and
we therefore score only functionality exposed through their Python
interfaces (for matUtils, via BTE).
Thus, underlying code may implement additional capabilities that are
not counted here if they are not available as documented, reusable
library operations.
For instance, ARGneedle-lib uses efficient tree iteration
to rapidly compute a GRM-vector product
using algorithms roughly analogous to tskit's \citep[described in][]{zhu2024fastvariancecomponent};
however, the tree iteration is only available in C++.
The comparisons were performed using tskit version 1.0.2,
arg-needle-lib version 1.2.1, BTE version 0.9.3, and DendroPy version 5.0.8.

\clearpage
\input{functionality_table.tex}

\bibliography{paper}

\end{document}
